\chapter{L'adversaire : trois stratégies mesurées}
\label{chap-adversaire}

\section{Une fonction pure}

L'adversaire artificiel se résume à une fonction (\cref{lst-ia}) : étant donné un niveau,
la connaissance publique du plateau visé et un générateur aléatoire, elle désigne la
prochaine case à viser.

\begin{lstlisting}[caption={L'adversaire tout entier (\fichier{include/battleship/ai.h}).},
                   label={lst-ia}]
bs_status bs_ai_choose_target(bs_difficulty difficulty,
                              const bs_knowledge *knowledge,
                              bs_rng *rng, bs_coord *target);
\end{lstlisting}

Cette fonction \textbf{ne garde aucune mémoire} d'un appel à l'autre. \enquote{Un bateau
est blessé} n'est pas une information à retenir : elle se lit sur le plateau, où une case
\code{HIT} désigne un pont touché d'un bateau encore à flot.

\begin{attention}
Dans la version~1, la poursuite d'un bateau touché reposait sur une variable globale. Après
un coup au but, l'ordinateur tirait au sort l'une des quatre voisines ; si elle était hors
du plateau ou déjà jouée, il \emph{abandonnait la poursuite} et retournait au tir aléatoire,
oubliant le bateau blessé. Pire, la case tirée au sort était lue dans le tableau
\emph{avant} que l'on vérifie qu'elle était sur le plateau : une lecture hors limites,
comportement indéfini en C.
\end{attention}

Une fonction sans mémoire ne peut pas se désynchroniser du jeu. Elle se prête aussi à toutes
les expériences : on peut lui soumettre n'importe quelle position, la faire jouer contre
elle-même, ou s'en servir pour conseiller le joueur.

\section{Trois niveaux, une compétence de plus à chaque marche}

\begin{table}[ht]
  \centering
  \begin{tabularx}{\textwidth}{@{}l>{\raggedright\arraybackslash}Xccc@{}}
    \toprule
    \textbf{Niveau} & \textbf{Stratégie} &
    \rotatebox{60}{\small ne rejoue pas} & \rotatebox{60}{\small déduit} &
    \rotatebox{60}{\small achève, anticipe} \\
    \midrule
    Facile    & Tire au hasard parmi les cases non encore visées. &
                $\bullet$ & & \\
    Moyen     & Tire au hasard, mais jamais sur une case que les règles prouvent vide. &
                $\bullet$ & $\bullet$ & \\
    Difficile & Achève les bateaux blessés ; en chasse, vise la case la plus probable. &
                $\bullet$ & $\bullet$ & $\bullet$ \\
    \bottomrule
  \end{tabularx}
  \caption{Les trois niveaux de l'adversaire.}
  \label{tab-niveaux}
\end{table}

Le niveau facile ignore les déductions du \cref{sec-deductions} : il peut tirer sur une
case marquée d'un drapeau. Le niveau moyen s'en sert et c'est sa seule différence. Le
niveau difficile y ajoute un raisonnement sur les bateaux qui restent.

\section{Le niveau difficile : une carte de densité}

\subsection{Principe}

Le niveau difficile se pose une question simple : \emph{parmi toutes les façons de placer
les bateaux encore à flot qui sont compatibles avec ce que je sais, combien passent par
cette case ?} La case où passent le plus de placements est la plus probable.

Notons $K$ la connaissance courante, $n_\ell$ le nombre de bateaux de longueur $\ell$ encore
à flot et $P_\ell(K)$ l'ensemble des positions cohérentes pour un tel bateau. Le score
d'une case $c$ encore inconnue est
\begin{equation}
  s(c) \;=\; \sum_{\ell = 1}^{4} n_\ell \cdot
  \bigl|\{\, p \in P_\ell(K) \;:\; c \in p \,\}\bigr| .
  \label{eq-densite}
\end{equation}

Une position $p$ est \emph{cohérente} avec $K$ lorsque les trois conditions suivantes sont
réunies :

\begin{enumerate}
  \item chacune de ses cases est soit inconnue, soit un pont touché d'un bateau à flot ;
  \item elle n'est pas faite \emph{uniquement} de ponts touchés --- ce bateau serait déjà
        coulé ;
  \item aucun pont touché ne la jouxte sans lui appartenir --- deux bateaux ne se touchent
        pas.
\end{enumerate}

La \cref{fig-densite} montre la carte obtenue sur un plateau vierge. Un coin n'est couvert
que par 16 placements, le centre par 36 : sans qu'aucune règle ne le lui dise, l'ordinateur
ouvre la partie au centre, là où les grands bateaux ont le plus de façons de se trouver.

\begin{figure}[ht]
  \centering
  \begin{tikzpicture}[scale=0.58]
    % Scores réels calculés par bs_ai_density() sur un plateau vierge (symétrique).
    \foreach \v [count=\i from 0] in {
      16,22,25,26,26,26,26,25,22,16,
      22,28,31,32,32,32,32,31,28,22,
      25,31,34,35,35,35,35,34,31,25,
      26,32,35,36,36,36,36,35,32,26,
      26,32,35,36,36,36,36,35,32,26,
      26,32,35,36,36,36,36,35,32,26,
      26,32,35,36,36,36,36,35,32,26,
      25,31,34,35,35,35,35,34,31,25,
      22,28,31,32,32,32,32,31,28,22,
      16,22,25,26,26,26,26,25,22,16} {
      \pgfmathtruncatemacro{\col}{mod(\i,10)}
      \pgfmathtruncatemacro{\row}{9 - int(\i/10)}
      \pgfmathtruncatemacro{\teinte}{(\v - 16) * 5}
      \fill[ocean!\teinte!ecume] (\col,\row) rectangle ++(1,1);
      \ifnum\v>29
        \node[font=\sffamily\scriptsize, text=white] at (\col+0.5,\row+0.5) {\v};
      \else
        \node[font=\sffamily\scriptsize, text=marine] at (\col+0.5,\row+0.5) {\v};
      \fi
    }
    \plateau{10}
    \foreach \l [count=\i from 0] in {A,B,C,D,E,F,G,H,I,J}
      \node[font=\sffamily\scriptsize, text=ardoise] at (\i+0.5,10.45) {\l};
    \foreach \n in {1,...,10}
      \node[font=\sffamily\scriptsize, text=ardoise] at (-0.45,10.5-\n) {\n};
  \end{tikzpicture}
  \caption{Carte de densité d'un plateau vierge : nombre de placements de la flotte
           couvrant chaque case (valeurs produites par \code{bs\_ai\_density()}).}
  \label{fig-densite}
\end{figure}

\subsection{Chasser et achever avec le même code}

Tant qu'aucun bateau n'est blessé, tous les placements cohérents comptent : c'est la
\emph{chasse}. Dès qu'une case \code{HIT} existe, on ne retient que les placements qui
\emph{expliquent} au moins un pont touché, pondérés par le nombre de ponts expliqués. La
même formule devient alors une procédure d'\emph{achèvement} :

\begin{itemize}
  \item autour d'un pont touché isolé, les quatre voisines reçoivent un score, d'autant
        plus élevé que de longs bateaux peuvent s'étendre dans leur direction ;
  \item dès que deux ponts alignés sont touchés, la condition~3 élimine toute position
        perpendiculaire : seules les deux cases prolongeant la ligne conservent un score.
\end{itemize}

Aucun cas particulier n'est programmé pour cela ; le comportement découle des trois
conditions de cohérence (\cref{alg-densite}).

\begin{algorithm}[ht]
  \caption{Choix d'une cible au niveau difficile}
  \label{alg-densite}
  \KwIn{la connaissance $K$, un générateur aléatoire $g$}
  \KwOut{une case à viser}
  $s(c) \gets 0$ pour toute case $c$\;
  $\mathit{achever} \gets$ il existe une case \textsc{touchée} dans $K$\;
  \ForEach{longueur $\ell$ telle que $n_\ell > 0$}{
    \ForEach{position $p$ d'un bateau de longueur $\ell$ sur le plateau}{
      \lIf{$p$ n'est pas cohérente avec $K$}{passer à la suivante}
      $h \gets$ nombre de cases \textsc{touchées} de $p$\;
      \lIf{$\mathit{achever}$ et $h = 0$}{passer à la suivante}
      $w \gets n_\ell \times h$ si $\mathit{achever}$, $n_\ell$ sinon\;
      \lForEach{case inconnue $c$ de $p$}{$s(c) \gets s(c) + w$}
    }
  }
  \Return une case de score maximal, tirée au sort avec $g$ en cas d'égalité\;
\end{algorithm}

Le coût est dérisoire : quatre longueurs, deux orientations, cent origines, des bateaux de
quatre cases au plus --- quelques milliers d'opérations par tir. Le tirage au sort entre
cases de même score évite que l'ordinateur joue toujours la même ouverture.

\section{Mesurer plutôt que supposer}

\subsection{Protocole}

Le programme \code{battleship-cli simulate} fait jouer les stratégies sans interface. Deux
expériences sont menées, entièrement déterminées par une graine :

\begin{description}[style=nextline, leftmargin=1.2em]
  \item[En solitaire] Une stratégie tire sur une flotte aléatoire jusqu'à la couler. On
    compte les tirs : il en faut au moins 20, au plus 100. Cette mesure isole la qualité de
    la stratégie.
  \item[En duel] Deux stratégies s'affrontent avec la règle réelle \enquote{touché, on
    rejoue}. Chaque appariement est joué en alternant le camp qui commence, pour annuler
    l'avantage du premier tir.
\end{description}

Tous les chiffres de ce chapitre proviennent de la commande
\code{battleship-cli simulate -{}-games 10000 -{}-seed 1} et peuvent être reproduits à
l'identique.

\subsection{Résultats}

Les niveaux sont nettement ordonnés (\cref{tab-solo,fig-histogramme}) : il faut en moyenne
\num{96.2} tirs au niveau facile, \num{65.8} au niveau moyen et \num{55.3} au niveau
difficile. En duel (\cref{tab-duels}), le niveau difficile l'emporte sur le moyen dans
\SI{87.5}{\percent} des parties, et le niveau moyen remporte \num{9999} de ses
\num{10000} parties contre le facile. La diagonale, proche de \SI{50}{\percent}, confirme
que l'alternance des camps neutralise bien l'avantage du premier tir.

\begin{table}[H]
  \centering
  \begin{tabular}{@{}l S[table-format=2.2] S[table-format=1.2] S[table-format=2.0]
                    S[table-format=2.0] S[table-format=3.0] S[table-format=3.0]@{}}
    \toprule
    \textbf{Niveau} & {\textbf{Moyenne}} & {\textbf{Écart-type}} & {\textbf{Min.}} &
    {\textbf{Médiane}} & {\textbf{9\textsuperscript{e} décile}} & {\textbf{Max.}} \\
    \midrule
    Facile    & 96.19 & 4.10 & 69 & 97 & 100 & 100 \\
    Moyen     & 65.83 & 6.07 & 44 & 66 &  74 &  85 \\
    Difficile & 55.27 & 6.80 & 28 & 56 &  64 &  79 \\
    \bottomrule
  \end{tabular}
  \caption{Nombre de tirs nécessaires pour couler une flotte aléatoire
           (\num{10000} parties par niveau).}
  \label{tab-solo}
\end{table}

\begin{figure}[H]
  \centering
  \begin{tikzpicture}
    \begin{axis}[
      width=0.96\textwidth, height=6.4cm,
      xlabel={Nombre de tirs pour couler la flotte},
      ylabel={Parties (sur \num{10000})},
      xmin=25, xmax=101, ymin=0,
      xtick={30,40,...,100}, ytick distance=500,
      scaled y ticks=false,
      yticklabel={\pgfmathprintnumber[fixed, 1000 sep={\,}]{\tick}},
      axis lines=left, axis line style={draw=ardoise},
      tick style={draw=ardoise}, tick label style={font=\sffamily\footnotesize},
      label style={font=\sffamily\small},
      ymajorgrids, grid style={draw=ardoise!20},
      legend style={draw=none, fill=none, font=\sffamily\small, at={(0.03,0.97)},
                    anchor=north west, legend cell align=left},
      const plot mark mid, area style, enlarge x limits=false,
    ]
      \addplot[draw=vert, fill=vert, fill opacity=0.30, line width=0.7pt]
        table[x=shots, y=hard] {data/benchmark.dat} \closedcycle;
      \addlegendentry{Difficile}
      \addplot[draw=ocean, fill=ocean, fill opacity=0.30, line width=0.7pt]
        table[x=shots, y=medium] {data/benchmark.dat} \closedcycle;
      \addlegendentry{Moyen}
      \addplot[draw=corail, fill=corail, fill opacity=0.30, line width=0.7pt]
        table[x=shots, y=easy] {data/benchmark.dat} \closedcycle;
      \addlegendentry{Facile}
    \end{axis}
  \end{tikzpicture}
  \caption{Distribution du nombre de tirs par niveau. Le niveau facile bute sur le
           plafond de 100 tirs ; les deux autres forment des cloches bien séparées.}
  \label{fig-histogramme}
\end{figure}

\begin{table}[H]
  \centering
  \begin{tabular}{@{}l S[table-format=3.1] S[table-format=3.1] S[table-format=2.1]@{}}
    \toprule
    & {\textbf{contre Facile}} & {\textbf{contre Moyen}} & {\textbf{contre Difficile}} \\
    \midrule
    \textbf{Facile}    &  49.6 &   0.0 &  0.0 \\
    \textbf{Moyen}     & 100.0 &  49.7 & 12.0 \\
    \textbf{Difficile} & 100.0 &  87.5 & 50.0 \\
    \bottomrule
  \end{tabular}
  \caption{Duels : pourcentage de parties gagnées par le niveau de la ligne
           (\num{10000} parties par case).}
  \label{tab-duels}
\end{table}

\subsection{Comment les niveaux ont été choisis}

Les trois niveaux ne résultent pas d'une intuition mais d'une comparaison. Quatre stratégies
candidates ont été mesurées pendant le développement (\cref{tab-candidats}).

\begin{table}[H]
  \centering
  \begin{tabularx}{\textwidth}{@{}>{\raggedright\arraybackslash}X S[table-format=2.1]
                               S[table-format=+2.1] l@{}}
    \toprule
    \textbf{Stratégie} & {\textbf{Tirs}} & {\textbf{Gain}} & \textbf{Retenue comme} \\
    \midrule
    A. Hasard pur & 96.2 & & Facile \\
    B. A + respect des déductions & 65.8 & -30.4 & Moyen \\
    C. B + poursuite classique des bateaux touchés & 58.2 & -7.6 & --- \\
    D. B + carte de densité (chasse et achèvement) & 55.3 & -2.9 & Difficile \\
    \bottomrule
  \end{tabularx}
  \caption{Stratégies candidates. Le gain est exprimé par rapport à la ligne précédente ;
           la stratégie~C a été mesurée sur \num{2000} parties.}
  \label{tab-candidats}
\end{table}

L'enseignement de ce tableau est contre-intuitif. Avec la règle de non-contact, la
compétence qui compte le plus n'est ni la poursuite des bateaux touchés ni le calcul de
probabilités : c'est le simple fait de \textbf{ne pas tirer là où les règles prouvent qu'il
n'y a rien}. À elle seule, cette discipline économise trente tirs. La poursuite classique en
ajoute moins de huit, et la carte de densité trois de plus.

Les niveaux ont donc été espacés d'après ces mesures : la stratégie~C, trop proche de~D,
n'aurait pas constitué un niveau intermédiaire perceptible.

\begin{retenir}[Un plancher difficile à abaisser]
Pourquoi le meilleur niveau a-t-il encore besoin de 55 tirs pour 20 cases à toucher ? À
cause des quatre bateaux d'une seule case. Aucun raisonnement ne permet de les
\emph{localiser} : une fois les grands bateaux coulés, il faut explorer presque toutes les
cases restantes. La carte de densité se rapproche de ce plancher ; elle ne peut pas le
supprimer.
\end{retenir}
